%LTeX: language=pt-BR
\section{Introdução e Justificativa}

A regularização $L_0$ e suas generalizações para esparsidade de grupos são ferramentas fundamentais em aprendizado de máquina e processamento de sinais, permitindo a seleção de variáveis e a interpretabilidade de modelos matemáticos complexos. Devido à natureza não-convexa e combinatória da penalidade $L_0$, a resolução eficiente desses problemas em larga escala permanece um desafio em aberto na otimização contínua. Outras regularizações, como híbridos de $L_1$ (Lasso) e $L_2$ (ridge) com $L_0$, também são de interesse recente por seu potencial para evitar overfitting, mantendo os benefícios da regularização pura $L_0$ \cite{fastselect}.

Durante o desenvolvimento do doutorado no Brasil, o foco central tem sido o desenvolvimento de heurísticas eficientes de primeira ordem para essas classes de problemas. Especificamente, desenvolvemos algoritmos baseados em Nonmonotone Spectral Proximal Gradient (NSPG) combinados com Descida de Coordenadas (CD) e métodos de busca local combinatória. Esses métodos não monótonos atuam de forma eficiente como um pré-processamento (identificação grosseira de suporte) para algoritmos do estado da arte (i.e., PGCCD \cite{fastselect}), superando desafios de dimensionalidade em ambientes de alta complexidade e correlação de paramêtros. A implementação foi desenvolvida utilizando a linguagem Julia, que permite alta performace e flexibilidade no desenvolvimento de algoritmos de otimização.

Apesar da alta eficiência computacional dos métodos de primeira ordem (espectrais), sua capacidade de atuar como pré-processadores globais e identificar boas regiões do espaço de busca poderia se beneficiar substancialmente de abordagens de segunda ordem. Aprimorar essa etapa com operadores proximais estruturados permite que o algoritmo se aproxime muito mais da variedade correta, garantindo que os métodos subsequentes de descida de coordenadas e algoritmos combinatórios tenham mais facilidade em realizar o refinamento local e encontrar soluções globalmente melhores. É neste contexto que a colaboração internacional proposta se fundamenta.

\subsection{Justificativa para a Escolha do Grupo Anfitrião}
O grupo de pesquisa liderado pelo Prof. Dominique Orban na Polytechnique Montréal (GERAD) possui liderança e excelência internacionalmente reconhecidas no desenvolvimento de métodos escaláveis de otimização não-suave e de segunda ordem. Trabalhos recentes do grupo, como o método de Levenberg-Marquardt para mínimos quadrados regularizados não-suaves \cite{aravkin2024levenberg} e métodos de regularização cúbica adaptativa escaláveis \cite{dussault2024scalable}, representam o estado da arte na resolução robusta de problemas complexos de otimização. Adicionalmente, análises recentes sobre a complexidade de métodos de região de confiança para otimização não-suave \cite{leconte2023complexity} oferecem uma base teórica sólida para extensões em regularizações descontínuas de interesse. Além disso, o grupo tem demonstrado forte interesse no estudo do comportamento algorítmico não monótono e no desenvolvimento de operadores proximais alternativos \cite{Leconte_2024}, o que se alinha perfeitamente com a trajetória de pesquisa atual do candidato.

Além do alinhamento teórico, o grupo hospedeiro é desenvolvedor principal do \texttt{JuliaSmoothOptimizers} (JSO), um ecossistema amplamente utilizado para otimização em Julia. A imersão neste ambiente de desenvolvimento propiciará uma sinergia única, permitindo não apenas o avanço teórico da pesquisa iniciada no Brasil, mas também sua conversão em software de alto impacto e uso global. Dessa forma, o estágio no exterior trará uma contribuição substancial que não poderia ser plenamente obtida na instituição de origem.
